\section{Introduction}
A telemetry monitor acts on the information that survives acquisition, retention,
and transmission. Evaluating a detector on a complete archive does not answer
which measurements a recorder should preserve before an engineer knows that a
fault has occurred. Onboard processing, storage, and compression between ground
contacts are established responsibilities of spacecraft data
systems~\citep{esa-data-systems}. These constraints motivate evaluating the retention
decision as well as the downstream detector.

Consider a short excursion followed by a return to the ordinary range. Keeping
only the final observation of a reporting cell can erase the excursion. Extrema
preserve magnitude but not necessarily exact timing, sequence, or cross-channel
relationships. Conversely, recording every unusual interval can exhaust a budget
on harmless changes. A useful recorder needs selective detail and a way to preserve
observations outside its trigger's current understanding of abnormal behavior.

We ask: \emph{under an explicit byte cap and bounded recorder state, which causal
retention choices preserve evidence useful to downstream receivers?} Here,
``fault evidence'' means measurements and temporal context supporting subsequent
assessment, not a validated root-cause diagnosis. A receiver may differ from the
capture trigger, and its input requirements may change which retained observations
are useful. Two isolated samples are not interchangeable with a consecutive pair
needed by a forecasting model.

This report makes three contributions:
\begin{enumerate}
\item An explicit application-byte contract and decoded-packet evaluation protocol
that preserve original timestamps and values, separate capture from labels, and
keep source-eligible event denominators fixed across policies.
\item A reproducible Mission1 development study with common byte caps, seeded
sampling controls, component ablations, capture/receiver mismatch, and illustrative
finite-buffer delivery. Both evidence and actual expenditure are reported.
\item A measured distinction between pointwise and temporal evidence preservation,
with an exploratory adjacent-pair control that partially addresses the observed
failure of singleton background capture.
\end{enumerate}

The analysis addresses three questions: how context changes when event counts do
not; what survives when the capture trigger misses an event; and how receiver input
requirements and delivery affect the evidence available for assessment. We do not
claim that triggered buffering or anomaly-guided prioritization is new. All mission
results concern the same 11 validation events, and the pair control was designed
after observing those results. Independent-mission and operator validation remain
necessary before claiming generalization or diagnostic benefit.
